\chapter*{Five complaints that changed the project}
\addcontentsline{toc}{chapter}{Five complaints that changed the project}

On 18 August 2026, a technically strong survey failed its first encounter with
a reader. The source ran to 2,326 lines. It reconstructed the mathematics of
Hamiltonian Monte Carlo for models with a zero lower bound, cited the relevant
papers, and documented the implementation in unusual detail. A mechanical
conversion had also made a mess of the typesetting. The reader's first verdict
was blunt: the mathematics did not look like mathematics.

Repairing the conversion did not solve the problem. It merely exposed the next
one. The survey read like an internal audit: file paths stood in for
explanations, inspection dates stood in for scope, and project labels appeared
before the objects they named. Once that register was removed, the reader
noticed the repeated sentence and section shapes. Once the repetition was
broken, defensive qualifications became conspicuous. Once the author replaced
those qualifications with reasons, the reader found the important omission:
the survey explained how to validate an algorithm but supplied no reference
models on which to do it.

Five rounds of feedback produced five different diagnoses. Figure
\ref{fig:zlb-rescue} shows their order. The order matters because each defect
had hidden the next one.

\hvFigureZLBRescue

The final manuscript ran to 49 pages, with 109 equations and 55 references.
The owner accepted it after the fifth repair. Throughout the rewrite, the
equation bodies remained byte-identical, labels and citations were counted,
and the rendered PDF was checked for the internal language that source-level
searches could miss. This is a useful engineering record. It is not an efficacy
study: there was one document, one owner, no control group, and no independent
reader panel.

\section*{The product hidden inside the rescue}

The rescue worked because four kinds of work were joined. A check located an
observable problem. An author decided what the passage needed to mean. A
comparison protected the equations, numbers, citations, and claims while the
prose changed. A reader then judged the rendered document without relying on
the repository history. Existing tools perform pieces of this sequence, but
the pieces do not share a record or end in a reader outcome.

Humanvoice is the proposed connective layer. It would sit beside the author's
editor and repository, not replace them. It would turn a reader brief into an
argument blueprint, draft in bounded units, and run a fail-closed pre-human
check before it turns a version into a reader packet. A broad complaint such as
``this is hard to read'' would become a small set of located questions; the
objects that carry technical meaning would remain protected; and the author's
decision would stay visible. It would never certify that prose is ``human,''
infer authorship from one document, silently rewrite a source, or ask an
expensive reader to serve as its ordinary debugger.

The wager is economic as much as editorial. A locally fluent document can
still consume several rounds of expert diagnosis and rereading. Humanvoice is
worth building only if it reduces that burden without increasing reader time,
false alarms, or meaning errors. A handsome score from a model is not the
outcome. A reader who can recover the problem, mechanism, evidence,
qualification, and requested action is.

\section*{The proposition under test}

The first investment is deliberately bounded: twelve weeks, a product
engineer, a research and evaluation lead, a domain editor, a part-time
security or policy owner, independent coding and reader review, a small
adjudicated corpus, and one live-document feasibility run. Weeks 1--6 deliver
the protected LaTeX core and its second-operator replay. The authoring
extension, model critics, and repair loop proceed only if that checkpoint
passes; otherwise the sponsor receives the protected review utility rather
than an unfinished platform.

At the end of twelve weeks the sponsor should be able to answer three
questions:

\begin{enumerate}
\item Can the system locate useful editorial questions and protect
      meaning-bearing objects on a real manuscript?
\item Can authors and readers use the complete path without the review burden
      exceeding the burden it is meant to remove?
\item Is a comparative study large enough to estimate an effect worth its
      expected cost?
\end{enumerate}

One live case can answer the first two as feasibility questions. It cannot
show that Humanvoice generally reduces revision rounds. That stronger claim
requires new documents, readers who did not tune the rules, a declared
comparator, and uncertainty reported around the result.

\begin{decisionbox}{Recommendation}
Fund the bounded build and feasibility run. Do not authorize deployment or an
efficacy claim. Continue only if the system protects meaning, keeps the reader
outcome observable, and produces a credible, priced design for the comparison
that follows.
\end{decisionbox}

\section*{Evidence strong enough to use today}

The research record makes the design plausible, but it is not complete.
Experiments show that generative assistance can save time on some writing and
knowledge tasks, while other studies show errors outside a model's capability
frontier, uneven gains across workers, convergence in assisted writing, and
bias in model-based judging. Research on writing and technical communication
explains why readers need visible actors, relations, evidence, and rhetorical
purpose. Software surveys identify mature components for parsing, linting,
comparison, citation identity, and local execution. None of those findings
establishes the effect of the combined Humanvoice workflow on expert technical
review.

The current audit passes broad retrieval, bibliographic identity checks, and a
bounded specialist import. Its harder challenge queries recover fourteen of
eighteen known items and therefore do not pass. Independent screening,
full-text appraisal of every load-bearing claim, held-out package benchmarks,
and external review also remain unfinished. The manuscript therefore
distinguishes source results from product inferences and treats completion of
those tasks as work, not as a confidence adjective.

\begin{table}[H]
\centering
\small
\begin{tabularx}{\textwidth}{@{}p{0.22\textwidth}L p{0.25\textwidth}@{}}
\toprule
Question & What the present record supports & What twelve weeks must add \\
\midrule
Is the failure real? & Several internal documents record correct work that
still required repeated reader-led reconstruction. & Baseline burden measured
on live documents outside the rule-writing set. \\
Can the pieces be built? & Mature components exist for most lower-level jobs;
the integration is technically plausible. & Held-out fixtures, explicit
fallbacks, and one complete local-first run. \\
Will the workflow help? & Adjacent studies justify the mechanism and the
outcomes chosen for measurement. & A later comparative study; one feasibility
case cannot estimate efficacy. \\
Can it be trusted? & Protected comparison, provenance, and named human review
give a testable safety model. & Observed false alarms, meaning errors, privacy
exceptions, and reader burden. \\
Is it worth paying for? & Saved expert attention is the relevant benefit. &
Local volume, loaded costs, baseline rounds, reader time, and maintenance
cost. \\
\bottomrule
\end{tabularx}
\caption{The proposal separates what is known from what the first build must
observe.}
\label{tab:orientation-route}
\end{table}

\section*{How the argument unfolds}

Part I starts with the reader's cost and then examines the case record that
made the failure visible. Part II asks what the scholarly and software
literatures actually establish, including the evidence against tempting
shortcuts such as automatic rewriting and document-level AI detection. Part
III follows a single teaching case through the proposed product, then states
the architecture and the deliberately narrow first build. Part IV specifies
the measurement, schedule, economics, adoption path, and stopping rules. The
appendices retain the schemas, fixtures, evidence ledger, and calculations
needed to reproduce the choices without forcing those records into the main
narrative.

The document can therefore be read at two depths. The chapters carry the
argument for a funding decision. The appendices let an implementer or reviewer
inspect how each consequential claim becomes a record, a test, or an open
question. Both belong in this volume because the proposal has to persuade a
reader and survive implementation.
